\documentclass[11pt,a4paper]{article}
% preamble.tex — shared setup for human-verification documents
% Compile each document with:  tectonic <File>.tex

\usepackage[T1]{fontenc}
\usepackage[utf8]{inputenc}
\usepackage{lmodern}
\usepackage{amssymb,amsmath,amsthm}
\usepackage{mathtools}
\usepackage[dvipsnames]{xcolor}
\usepackage{listings}
\usepackage{hyperref}
\usepackage[margin=2.5cm]{geometry}
\usepackage{parskip}

% ── Colours ──────────────────────────────────────────────────────────
\definecolor{leanBg}{HTML}{F7F7F0}
\definecolor{haskellBg}{HTML}{F0F4F7}
\definecolor{keyword}{HTML}{7B3294}
\definecolor{comment}{HTML}{5A8C32}
\definecolor{noteFg}{HTML}{B35900}

% ── Lean 4 listings language ─────────────────────────────────────────
% NOTE: Lean listing blocks use ASCII approximations of unicode symbols
% (e.g. -> for arrow, o. for compose).  The literate key substitutes
% these with proper math-mode renderings in the PDF.
\lstdefinelanguage{Lean4}{
  morekeywords={def,theorem,structure,where,namespace,end,import,fun,
                let,have,show,by,exact,intro,cases,rfl,subst,obtain,
                funext,simp,rw,apply,variable,Type,Prop,instance,
                section,open,class,extends,deriving,sorry,noncomputable},
  sensitive=true,
  morecomment=[l]{--},
  morecomment=[n]{/-}{-/},
  morestring=[b]",
}

% ── Lean listing style ───────────────────────────────────────────────
\lstdefinestyle{lean}{
  language=Lean4,
  backgroundcolor=\color{leanBg},
  basicstyle=\small\ttfamily,
  keywordstyle=\color{keyword}\bfseries,
  commentstyle=\color{comment}\itshape,
  breaklines=true,
  frame=single,
  framerule=0.4pt,
  xleftmargin=4pt,
  xrightmargin=4pt,
  aboveskip=6pt,
  belowskip=4pt,
  literate=
    {->}{{$\to$}}2
    {=>}{{$\Rightarrow$}}2
    {<|}{{$\langle$}}1
    {|>}{{$\rangle$}}1
    {==>}{{$\Longrightarrow$}}3
    { o. }{{$\;\circ\;$}}3
    {forall}{{$\forall$}}1
    {:=}{{$\,:=\,$}}2
    {!=}{{$\neq$}}2
    {<=}{{$\leq$}}2
    {>=}{{$\geq$}}2,
}

% ── Haskell listing style ───────────────────────────────────────────
\lstdefinestyle{haskell}{
  language=Haskell,
  backgroundcolor=\color{haskellBg},
  basicstyle=\small\ttfamily,
  keywordstyle=\color{keyword}\bfseries,
  commentstyle=\color{comment}\itshape,
  breaklines=true,
  frame=single,
  framerule=0.4pt,
  xleftmargin=4pt,
  xrightmargin=4pt,
  aboveskip=6pt,
  belowskip=4pt,
  literate=
    {->}{{$\to$}}2
    {=>}{{$\Rightarrow$}}2,
}

% ── Theorem environments ────────────────────────────────────────────
\theoremstyle{definition}
\newtheorem{defn}{Definition}[section]
\newtheorem{thm}[defn]{Theorem}
\newtheorem{prop}[defn]{Proposition}

\theoremstyle{remark}
\newtheorem*{fidelitynote}{Fidelity note}

% ── Convenience macros ───────────────────────────────────────────────
\newcommand{\lean}[1]{\texttt{#1}}      % inline Lean identifier
\newcommand{\hask}[1]{\texttt{#1}}      % inline Haskell identifier
\newcommand{\Ext}{\mathrm{Ext}}
\newcommand{\Pos}{\mathrm{Pos}}
\newcommand{\Shape}{\mathrm{Shape}}

\usepackage{stmaryrd}

\title{Human Verification: \texttt{Composition}\\[4pt]
  \large Container Composition $\triangleleft$ and its Monoidal Structure}
\author{Machine-generated from \texttt{lean/Composition.lean} for review by Neil Ghani}
\date{July 2026}

\begin{document}
\maketitle

\begin{abstract}
This document presents every definition and theorem \emph{statement} from the
Lean~4 module \lean{Composition} in three parallel forms: the original Lean,
a readable Haskell rendering, and precise mathematical English.  Proofs are
omitted --- they are machine-checked by Lean.

The module equips the category $\mathbf{Cont}$ of containers with a
\textbf{composition product} $\triangleleft$ whose extension satisfies
$\llbracket G \triangleleft F \rrbracket = \llbracket G \rrbracket \circ
\llbracket F \rrbracket$.  The unit is the identity container $I$.  The main
results are the left and right unitors and the associator, each proved as an
explicit isomorphism of containers.  A $\triangleleft$-comonoid is exactly a
directed container.
\end{abstract}

\tableofcontents
\bigskip

%===================================================================
\section{Overview}
%===================================================================

For containers $F = S \triangleleft P$ and $G = T \triangleleft Q$, their
\textbf{composition} $G \triangleleft F$ (``$G$ outside, $F$ inside'') has
\begin{itemize}
  \item shapes $\Sigma_{t : T}(Q(t) \to S)$ --- an outer $G$-shape $t$ together
    with an assignment of an inner $F$-shape to each $G$-position; and
  \item positions at $(t, f)$ given by $\Sigma_{q : Q(t)} P(f(q))$ --- pick an
    outer position $q$, then a position in the inner shape $f(q)$.
\end{itemize}

The key fact is that $\llbracket G \triangleleft F \rrbracket X \cong \llbracket
G \rrbracket (\llbracket F \rrbracket X)$, made definitional by the
\lean{compExt}/\lean{compExtInv} pair.

%===================================================================
\section{The composition product}
%===================================================================

\subsection{Definition: \lean{Container.comp}}

\begin{defn}[Container composition]
For containers $G$ and $F$, their composition $G \triangleleft F$ is the
container with
\[
  \mathrm{Shape}(G \triangleleft F) = \Sigma_{t : G.\mathrm{Shape}}
    \bigl(G.\mathrm{Pos}(t) \to F.\mathrm{Shape}\bigr),
\]
\[
  \mathrm{Pos}(G \triangleleft F)(t, f) = \Sigma_{q : G.\mathrm{Pos}(t)}
    F.\mathrm{Pos}(f(q)).
\]
\end{defn}

\paragraph{Lean.}
\begin{lstlisting}[style=lean]
def Container.comp (G F : Container) : Container where
  Shape := (t : G.Shape) * (G.Pos t -> F.Shape)
  Pos := fun tf => (q : G.Pos tf.1) * F.Pos (tf.2 q)

scoped infixr:70 " comp " => Container.comp
-- (written with the unicode diamond operator in Lean source)
\end{lstlisting}

\paragraph{Haskell.}
\begin{lstlisting}[style=haskell]
-- G `outside` F:
data CompShape g p_g f = CompShape g (p_g -> f)
-- Positions at (t, assign):
data CompPos p_g p_f t assign = CompPos (p_g t) (p_f (assign (p_g t)))
-- Approximate (non-dependent):
data Comp g p_g f p_f x = Comp g (p_g -> f) (p_g -> p_f -> x)
\end{lstlisting}

\begin{fidelitynote}
The shape and position types depend on one another in a genuinely dependent
way ($P_F$ is indexed by the value of the inner shape function $f(q)$).
Haskell cannot express this; the rendering above is an approximation that
loses the dependency.
\end{fidelitynote}


\subsection{Definition: \lean{Container.I}}

\begin{defn}[Identity container]
The \textbf{identity container} $I$ has one shape (the unit type) and one
position.  Its extension is the identity functor: $\llbracket I \rrbracket X
= X$.
\end{defn}

\paragraph{Lean.}
\begin{lstlisting}[style=lean]
def Container.I : Container where
  Shape := Unit
  Pos := fun _ => Unit
\end{lstlisting}

\paragraph{Haskell.}
\begin{lstlisting}[style=haskell]
-- Identity functor: Id x = x.
-- As a container: one shape (Unit), one position (Unit).
newtype Id x = Id x
\end{lstlisting}

%===================================================================
\section{Extension isomorphism}
%===================================================================

\subsection{Definition: \lean{Container.compExt} and \lean{Container.compExtInv}}

\begin{defn}[Extension isomorphism]
There is a natural bijection
\[
  \llbracket G \triangleleft F \rrbracket X \;\cong\; \llbracket G \rrbracket
  \bigl(\llbracket F \rrbracket X\bigr),
\]
given by:
\begin{align*}
  \mathrm{compExt}\,\langle (t, f), \ell \rangle
    &= \langle t,\; q \mapsto \langle f(q),\; r \mapsto \ell(q, r) \rangle \rangle, \\
  \mathrm{compExtInv}\,\langle t, m \rangle
    &= \langle (t,\; q \mapsto m(q).1),\; (q, r) \mapsto m(q).2(r) \rangle.
\end{align*}
\end{defn}

\paragraph{Lean.}
\begin{lstlisting}[style=lean]
def Container.compExt {G F : Container} {X : Type} :
    Ext (G comp F) X -> Ext G (Ext F X) :=
  fun p => <|p.1.1, fun q => <|p.1.2 q, fun r => p.2 <|q, r|>|>|>

def Container.compExtInv {G F : Container} {X : Type} :
    Ext G (Ext F X) -> Ext (G comp F) X :=
  fun p => <|<|p.1, fun q => (p.2 q).1|>, fun qr => (p.2 qr.1).2 qr.2|>
\end{lstlisting}

\paragraph{Haskell.}
\begin{lstlisting}[style=haskell]
compExt :: Comp g p_g f p_f x -> g -> (p_g -> (f, p_f -> x))
compExt (Comp t assign lab) = (t, \q -> (assign q, \r -> lab q r))

compExtInv :: (g, p_g -> (f, p_f -> x)) -> Comp g p_g f p_f x
compExtInv (t, m) = Comp t (fst . m) (\q r -> snd (m q) r)
\end{lstlisting}


\subsection{Theorems: \lean{compExtInv\_compExt} and \lean{compExt\_compExtInv}}

\begin{thm}[Extension isomorphism is an isomorphism]
The maps \lean{compExt} and \lean{compExtInv} are mutually inverse:
\begin{align*}
  \mathrm{compExtInv} \circ \mathrm{compExt} &= \mathrm{id}, \\
  \mathrm{compExt} \circ \mathrm{compExtInv} &= \mathrm{id}.
\end{align*}
Both round trips hold definitionally (proved by \lean{rfl} after
\lean{cases}), because the position fibres agree definitionally.
\end{thm}

\paragraph{Lean.}
\begin{lstlisting}[style=lean]
theorem Container.compExtInv_compExt {G F : Container} {X : Type}
    (p : Ext (G comp F) X) :
    Container.compExtInv (Container.compExt p) = p

theorem Container.compExt_compExtInv {G F : Container} {X : Type}
    (p : Ext G (Ext F X)) :
    Container.compExt (Container.compExtInv p) = p
\end{lstlisting}

\paragraph{Haskell.}
\begin{lstlisting}[style=haskell]
-- compExtInv . compExt = id   (and vice versa)
-- Both hold by the eta-laws for pairs.
\end{lstlisting}

%===================================================================
\section{Container isomorphisms}
%===================================================================

\subsection{Structure: \lean{ContainerIso}}

\begin{defn}[Container isomorphism]
An \textbf{isomorphism of containers} $C \cong D$ is a pair of container
morphisms $\varphi : C \to D$ and $\psi : D \to C$ satisfying:
\[
  \varphi \circ \psi = \mathrm{id}_C \quad\text{and}\quad
  \psi \circ \varphi = \mathrm{id}_D.
\]
(Composition here is morphism composition in $\mathbf{Cont}$.)
\end{defn}

\paragraph{Lean.}
\begin{lstlisting}[style=lean]
structure ContainerIso (C D : Container) where
  hom : ContainerMorphism C D
  inv : ContainerMorphism D C
  hom_inv : hom.comp inv = ContainerMorphism.id C
  inv_hom : inv.comp hom = ContainerMorphism.id D
\end{lstlisting}

\paragraph{Haskell.}
\begin{lstlisting}[style=haskell]
data ContainerIso c d = ContainerIso
  { hom     :: Morph c d
  , inv     :: Morph d c
  , hom_inv :: hom `compMorph` inv == idMorph c   -- laws
  , inv_hom :: inv `compMorph` hom == idMorph d
  }
\end{lstlisting}

%===================================================================
\section{Monoidal coherences}
%===================================================================

\subsection{Definition: \lean{Container.leftUnitor}}

\begin{defn}[Left unitor]
There is a canonical isomorphism of containers
\[
  \lambda_F : I \triangleleft F \;\cong\; F.
\]
A shape of $I \triangleleft F$ is $(\langle\rangle, f)$ with
$f : \mathrm{Unit} \to F.\mathrm{Shape}$; the bijection with $F.\mathrm{Shape}$
sends $f \mapsto f(\langle\rangle)$ and $s \mapsto (\lambda\, \_ \to s)$.
Positions collapse by the same $\mathrm{Unit}$-projection.
\end{defn}

\paragraph{Lean.}
\begin{lstlisting}[style=lean]
def Container.leftUnitor (F : Container) :
    ContainerIso (Container.I comp F) F where
  hom :=
    { onShapes := fun tf => tf.2 ()
      onPos := fun _ p => <|(), p|> }
  inv :=
    { onShapes := fun s => <|(), fun _ => s|>
      onPos := fun _ qp => qp.2 }
  hom_inv := rfl
  inv_hom := rfl
\end{lstlisting}

\paragraph{Haskell.}
\begin{lstlisting}[style=haskell]
leftUnitor :: ContainerIso (I `Comp` F) F
-- hom: ((), f) |-> f (); pos: p |-> ((), p)
-- inv: s     |-> ((), const s); pos: (q,p) |-> p
\end{lstlisting}


\subsection{Definition: \lean{Container.rightUnitor}}

\begin{defn}[Right unitor]
There is a canonical isomorphism of containers
\[
  \rho_F : F \triangleleft I \;\cong\; F.
\]
A shape of $F \triangleleft I$ is $(s, g)$ with $g : F.\mathrm{Pos}(s) \to
\mathrm{Unit}$; the bijection projects onto $s$.  Positions collapse
$\Sigma_{p : F.\mathrm{Pos}(s)} \mathrm{Unit} \cong F.\mathrm{Pos}(s)$.
\end{defn}

\paragraph{Lean.}
\begin{lstlisting}[style=lean]
def Container.rightUnitor (F : Container) :
    ContainerIso (F comp Container.I) F where
  hom :=
    { onShapes := fun sg => sg.1
      onPos := fun _ p => <|p, ()|> }
  inv :=
    { onShapes := fun s => <|s, fun _ => ()|>
      onPos := fun _ pu => pu.1 }
  hom_inv := rfl
  inv_hom := rfl
\end{lstlisting}

\paragraph{Haskell.}
\begin{lstlisting}[style=haskell]
rightUnitor :: ContainerIso (F `Comp` I) F
-- hom: (s, g) |-> s;   pos: p |-> (p, ())
-- inv: s      |-> (s, const ()); pos: (p,()) |-> p
\end{lstlisting}


\subsection{Definition: \lean{Container.associator}}

\begin{defn}[Associator]
There is a canonical isomorphism of containers
\[
  \alpha_{H,G,F} : (H \triangleleft G) \triangleleft F \;\cong\;
  H \triangleleft (G \triangleleft F).
\]
On shapes it curries/uncurries the inner shape assignment:
a shape $((a, g), k)$ on the left (where $k : \Sigma_{u : H.\mathrm{Pos}(a)}
G.\mathrm{Pos}(g(u)) \to F.\mathrm{Shape}$) corresponds to $(a, m)$ on the
right (where $m(u) = \langle g(u),\, v \mapsto k(u, v) \rangle$).
On positions it re-associates the nested dependent sum:
$(( u, v ), p) \leftrightarrow (u, (v, p))$.

Crucially, the position fibre $F.\mathrm{Pos}((m\,u).2\,v)$ equals
$F.\mathrm{Pos}(k(u, v))$ \emph{definitionally}, so no transport is needed.
\end{defn}

\paragraph{Lean.}
\begin{lstlisting}[style=lean]
def Container.associator (H G F : Container) :
    ContainerIso ((H comp G) comp F) (H comp (G comp F)) where
  hom :=
    { onShapes := fun xk =>
        <|xk.1.1, fun u => <|xk.1.2 u, fun v => xk.2 <|u, v|>|>|>
      onPos := fun _ uvp => <|<|uvp.1, uvp.2.1|>, uvp.2.2|> }
  inv :=
    { onShapes := fun am =>
        <|<|am.1, fun u => (am.2 u).1|>, fun y => (am.2 y.1).2 y.2|>
      onPos := fun _ yp => <|yp.1.1, yp.1.2, yp.2|> }
  hom_inv := ...  -- ContainerMorphism.ext' rfl (cases + rfl)
  inv_hom := ...  -- ContainerMorphism.ext' rfl (cases + rfl)
\end{lstlisting}

\paragraph{Haskell.}
\begin{lstlisting}[style=haskell]
associator :: ContainerIso ((H `Comp` G) `Comp` F) (H `Comp` (G `Comp` F))
-- hom shapes: ((a,g),k) |-> (a, \u -> (g u, \v -> k (u,v)))
-- hom pos:    ((u,v),p) |-> (u,(v,p))
-- inv shapes: (a,m)     |-> ((a, fst . m), \(u,v) -> snd (m u) v)
-- inv pos:    (u,(v,p)) |-> ((u,v),p)
\end{lstlisting}

%===================================================================
\section{The compositional payoff}
%===================================================================

The fact that $({\mathbf{Cont}}, \triangleleft, I)$ is a monoidal category
(with the above coherences) means that a \textbf{$\triangleleft$-comonoid} in
$\mathbf{Cont}$ --- a container $C$ equipped with a counit $C \to I$ and a
comultiplication $C \to C \triangleleft C$ satisfying the comonoid laws ---
is exactly a \textbf{directed container} (equivalently, a comonad on
$\mathbf{Set}$ whose underlying functor is a container).  The comonoid laws
unfold to the five directed-container axioms (D1--D5) formalised in
\lean{Containers.Directed}.

%===================================================================
\section{Translation conventions}
%===================================================================

\begin{enumerate}
\item \textbf{Dependent types.}  The position type in $G \triangleleft F$
  depends on the shape value.  Haskell approximations drop this dependency.

\item \textbf{Definitional equalities.}  The round trips for \lean{compExt} /
  \lean{compExtInv} and for the associator both hold \emph{definitionally}
  in Lean (proved by \lean{rfl} after \lean{cases}).  This is a key feature
  of the Lean 4 type theory, not available in most proof assistants.

\item \textbf{Unicode notation.}  Lean uses the diamond operator
  $\triangleleft$ for container composition.  The listings use
  \texttt{comp} as an ASCII placeholder.

\item \textbf{Proof bodies omitted.}  Only statements are presented here.
\end{enumerate}

\end{document}
